\documentclass[a4paper]{report}
\linespread{1.5} % ajouter un espacement entre les lignes de tout le document 
%====================== PACKAGES ======================
%\documentclass{article}
\usepackage{xcolor}
\usepackage{listings}
\definecolor{mGreen}{rgb}{0,0.6,0}
\definecolor{mGray}{rgb}{0.5,0.5,0.5}
\definecolor{mPurple}{rgb}{0.58,0,0.82}
\definecolor{backgroundColour}{rgb}{0.95,0.95,0.92}

\lstdefinestyle{CStyle}{
    backgroundcolor=\color{backgroundColour},   
    commentstyle=\color{mGreen},
    keywordstyle=\color{magenta},
    numberstyle=\tiny\color{mGray},
    stringstyle=\color{mPurple},
    basicstyle=\footnotesize,
    breakatwhitespace=false,         
    breaklines=true,                 
    captionpos=b,                    
    keepspaces=true,                 
    numbers=left,                    
    numbersep=5pt,                  
    showspaces=false,                
    showstringspaces=false,
    showtabs=false,                  
    tabsize=2,
    language=java
}







\usepackage{algorithm}
\usepackage{algorithmic}
\usepackage[french]{babel}
\usepackage[utf8x]{inputenc}
%pour gérer les positionnement d'images
\usepackage{float}
\usepackage{amsmath}
\usepackage{graphicx}
\usepackage[colorinlistoftodos]{todonotes}
\usepackage{url}
%pour les informations sur un document compilé en PDF et les liens externes / internes
\usepackage{hyperref}
%pour la mise en page des tableaux
\usepackage{array}
\usepackage{tabularx}
%pour utiliser \floatbarrier
%\usepackage{placeins}
%\usepackage{floatrow}
%espacement entre les lignes
\usepackage{setspace}
%modifier la mise en page de l'abstract
\usepackage{abstract}
%police et mise en page (marges) du document
\usepackage[T1]{fontenc}
\usepackage[top=2cm, bottom=2cm, left=2cm, right=2cm]{geometry}
%Pour les galerie d'images
\usepackage{subfig}

%====================== INFORMATION ET REGLES ======================

%rajouter les numérotation pour les \paragraphe et \subparagraphe
\setcounter{secnumdepth}{4}
\setcounter{tocdepth}{4}

\hypersetup{							% Information sur le document
pdfauthor = {EL MRABTI Hamza},			% Auteurs
pdftitle = {Système de santé},			% Titre du document
pdfsubject = {Système de santé},		% Sujet
pdfkeywords = {Tag1, Tag2, Tag3, ...},	% Mots-clefs
pdfstartview={FitH}}					% ajuste la page à la largueur de l'écran
%pdfcreator = {MikTeX},% Logiciel qui a crée le document
%pdfproducer = {}} % Société avec produit le logiciel

%======================== DEBUT DU DOCUMENT ========================

\begin{document}

%régler l'espacement entre les lignes
\newcommand{\HRule}{\rule{\linewidth}{0.5mm}}

%page de garde
\begin{titlepage}
\begin{center}

% Upper part of the page. The '~' is needed because only works if a paragraph has started.
\includegraphics[width=1\textwidth]{./ensiass}~\\[1.5cm]

\textsc{\Large Ecole Nationale Supérieure d’Informatique et d’Analyse des Systèmes - RABAT }\\[2cm]

\textsc{\Large Projet de fin de la deuxième années
 }\\[0.5cm]

% Title
\HRule \\[0.4cm]

{\huge \bfseries   
 Gestion des données d'un patient dans un système de santé muti-centres
 \\[0.4cm] }
\HRule \\
[2cm]
% Author and supervisor
\begin{minipage}{0.4\textwidth}
\begin{flushleft} \large
\emph{Réalisé par :}\\
Hamza \textsc{El Mrabti}\\
Yahya \textsc{Benzha}\\ [1cm]

\emph{Filière :} \\
\textsc{Génie Logiciel }\\



\end{flushleft}
\end{minipage}
\begin{minipage}{0.4\textwidth}
\begin{flushright} \large
\emph{Encadré par:} \\
Pr. Ahmed \textsc{Ettalbi}\\[1cm]
\emph{Memebres du jury:} \\
Pr. Ahmed \textsc{Ettalbi}\\
Pr. Mounia \textsc{Abik}\\[1cm]


\end{flushright}
\end{minipage}

\vfill

% Bottom of the page
{\large \ Année Scolaire 2021/2022}

\end{center}
\end{titlepage}

%page blanche
\newpage
~
%ne pas numéroter cette page
\thispagestyle{empty}
%\newpage  ==> pour créer une nouvelle page



\chapter*{Remerciements}



\begin{center}
``La reconnaissance est la plus belle fleur qui jaillit de l'âme "
\end{center}
\begin{center} $\sim$ Henry Beecher . \end{center}
Nous voudrions tout d'abord adresser toute notre gratitude au 
Pr. Ahmed Ettalbi, notre encadrant, pour sa patience, sa disponibilité, sa rigueur scientifique, ses qualités humaines
et surtout ses judicieux conseils, qui ont contribué à alimenter notre
réfexion.

Ensuite nous exprimons notre profonde reconnaissance aux membres
du jury, pour le temps précieux qu'ils accordent volontairement en contribuant au bon déroulement du module.

Enfin nous tenons à remercier chaleureusement nos familles, nos amis et toutes les personnes qui ont contribué énormément à la réalisation de notre projet et qui nous ont aidés lors de la rédaction de ce rapport, aussi bien pour leurs encouragements et soutiens dans les moments forts.


\chapter*{Résumé}

ici c'est le résumé


\chapter*{Abstract}


This report concerns the project for the end of the second year of the engineering cycle, software engineering sector at the Ecole Supérieure d'Informatique et d'Analyse des Systèmes

The subject of this project was the management of patient data in a multi-center system.

  Our work was done in several phases:
  An analysis and design phase and a development phase.
  Thanks to the help of our supervisor, we have learned, over the past few months, to use
different tools and discover several concepts that helped us to develop this application.
  We will therefore present to you throughout this report the steps we have
followed as well as the tools we used to carry out our project.

\textbf{Keywords} : Health system , Patient data ,  Medical records , e-Health 


\tableofcontents
\thispagestyle{empty}
\setcounter{page}{0}
\listoffigures
\thispagestyle{empty}


\chapter*{Introduction générale}


ici c'est l'introduction générale


\setcounter{page}{0}
%ne pas numéroter le sommaire

% \newpage

%espacement entre les lignes d'un tableau
\renewcommand{\arraystretch}{1.5}

%====================== INCLUSION DES PARTIES ======================

~
\thispagestyle{empty}
%recommencer la numérotation des pages à "1"
\setcounter{page}{0}
\newpage





\chapter{chapter} 

\section{section }

\begin{figure}[!h]
\begin{center}
%taille de l'image en largeur
%remplacer "width" par "height" pour régler la hauteur
\includegraphics[height=10cm]{ensias.png}
\end{center}
%légende de l'image
\caption{legende de l'image  }
\end{figure}



%\input{./existant.tex}

 %\input{./besoins.tex}


%\input{./resultats.tex}



\newpage

%récupérer les citation avec "/footnotemark"
\nocite{*}

%choix du style de la biblio
\bibliographystyle{plain}
%inclusion de la biblio
\bibliography{bibliographie.bib}
%voir wiki pour plus d'information sur la syntaxe des entrées d'une bibliographie
\thispagestyle{empty}
\end{document}